% Reconstructed text from V5 PDF page 42; non-native source
\par\smallskip\noindent\rule{\textwidth}{0.4pt}\par{\sffamily\footnotesize Original V5 page 42 -- recovered text}\par\begin{flushleft}\footnotesize\ttfamily
\par\smallskip
\par\smallskip
\noindent Diagnostics and Serving Infrastructure for MAS 42\par
\par\smallskip
\par\smallskip
\par\smallskip
\noindent gap varies across retention strategies and budgets (Table 10). S4 and S3 have zero seed variance\par
\noindent (deterministic strategies); only S2 has genuine uncertainty bars.\par
\par\smallskip
\noindent Recovery-ablation: confirmed mechanism. A controlled follow-up experiment isolated the\par
\noindent recovery-detection heuristic as the mechanism behind the inverted slope. The original causal-\par
\noindent graph heuristic marks an error as “recovered” when the failure-causing agent subsequently pro-\par
\noindent duces any non-error event, suppressing it from the attribution. We tested three conditions (Ta-\par
\noindent ble 11):\par
\par\smallskip
\noindent Table 11: Recovery-ablation results on 220 TraceElephant traces (n = 220; deterministic for C2/C3).\par
\noindent McNemar’s test for paired nominal data.\par
\par\smallskip
\noindent  Condition Accuracy McNemar p 95\% CI\par
\noindent  C1 Original heuristic, 100\% retention 0.459 (101/220) — —\par
\noindent  C2 No-recovery logic, 100\% retention 0.532 (117/220) 0.004 vs C1 [+2.73, +11.82] pp\par
\noindent  C3 No-recovery + S4 importance, 10\% 0.541 (119/220) 0.892 vs C2 crosses zero\par
\noindent  McNemar exact test with mid-p correction. C1→C2 discordant pairs: b=6 (C2 correct, C1 wrong), c=22 (C1\par
\noindent  correct, C2 wrong). C2→C3: b=26, c=28. CI expressed as C2−C1 (improvement direction, in percentage\par
\noindent  points).\par
\par\smallskip
\noindent  Correcting the recovery heuristic accounts for most of the apparent subsampling benefit:\par
\noindent C2 recovers +7.3 percentage points over C1 (McNemar p = 0.004), matching the accuracy of\par
\noindent importance-based subsampling at 10\% budget under the original heuristic (C2 ≈ S4@10\% =\par
\noindent 0.532). Two independent paths — (a) fixing the heuristic directly, (b) subsampling away the\par
\noindent events the flawed logic depends on — converge to the same accuracy. This confirms that the\par
\noindent inverted slope in Table 10 was driven by the recovery-detection artifact, not by a genuine benefit\par
\noindent of less data.\par
\noindent  The residual 0.9-percentage-point improvement under importance sampling (C3 vs. C2) is\par
\noindent not statistically significant (p = 0.892) and should not be interpreted as robust without replica-\par
\noindent tion with other attributors. The improvement is not uniform across suites: GAIA suites show\par
\noindent the largest gains (+9.5 pp) while SWE-Agent slightly regresses under C2 (−2.3 pp), suggesting\par
\noindent the recovery logic interacts differently with framework-specific trace structures.\par
\par\smallskip
\noindent Implications. This finding characterizes a specific heuristic–retention interaction with a con-\par
\noindent firmed and statistically tested mechanism. Different attribution methods (e.g., LLM-judge ap-\par
\noindent proaches such as CHIEF or TrajAudit) would likely produce different strategy rankings and\par
\noindent may not exhibit the inverted slope, since they do not rely on the “recovered error” heuristic.\par
\noindent The practical takeaway is that retention policy design must account for the attribution method’s\par
\noindent assumptions: a retention strategy that improves one attributor may be neutral or harmful for\par
\noindent another. Replication with step-level attribution and LLM-judge methods remains necessary to\par
\noindent establish generality.\par
\par\smallskip
\par\smallskip
\noindent 8 DRICE Evaluation Framework\par
\noindent Existing comparison tables in multi-agent diagnostics and serving reviews typically list systems\par
\noindent by method category and report headline accuracy or throughput numbers. This makes cross-\par
\noindent layer comparison difficult: a diagnostic method with 70\% attribution accuracy and zero replay\par
\noindent cost is not directly comparable to an infrastructure system with 3× throughput gain and 74\%\par
\end{flushleft}
